\chapter{Formes quadratiques}\label{ch:b2:quadratic}

Une \omterm{def:b2:quadratic:def}{forme quadratique} est l'ombre algébrique d'une géométrie~: les
parties de signature nulle aplatissent, les parties positives
courbent dans un sens, les négatives dans l'autre. Ce chapitre
réduit toute \omterm{def:b2:quadratic:def}{forme quadratique} réelle à une somme de carrés $\pm$
(Gauss), prouve que les nombres de signes sont intrinsèques
(Sylvester), puis couronne la géométrie euclidienne par le
\emph{théorème spectral}~: les \omterm{def:b2:quadratic:adjoint}{endomorphismes symétriques} se
diagonalisent dans des bases orthonormées --- le théorème le plus
utilisé de l'\omterm{def:b2:structures:algebra}{algèbre} linéaire appliquée.

\section{Formes bilinéaires et quadratiques}

\begin{definition}\label{def:b2:quadratic:def}
Une \emph{forme bilinéaire \omterm{def:b2:quadratic:adjoint}{symétrique}} sur un espace vectoriel réel
$E$ est une application bilinéaire $\varphi \colon E \times E \to
\R$ vérifiant $\varphi(x, y) = \varphi(y, x)$~; la \emph{forme
quadratique}\index{forme quadratique} associée est $q(x) =
\varphi(x, x)$. On retrouve $\varphi$ à partir de $q$ par
\emph{polarisation}\index{identité de polarisation}~:
\[
\varphi(x, y) = \tfrac12\bigl(q(x + y) - q(x) - q(y)\bigr).
\]
Dans une base $(e_i)$, la \emph{matrice} de $\varphi$ est la matrice
\omterm{def:b2:quadratic:adjoint}{symétrique} $B = (\varphi(e_i, e_j))$, avec $q(x) = X^{\mathsf T} B
X$~; un changement de base de matrice $P$ remplace $B$ par
$P^{\mathsf T} B P$ (\emph{congruence}\index{matrices congruentes}
--- et non similitude~!). Le \emph{rang} de $q$ est
$\operatorname{rk} B$ (invariant~: la congruence multiplie par des
matrices inversibles).
\end{definition}

\begin{example}
Sur $\R^2$~: $q(x, y) = x^2 + 4xy + y^2$ a pour matrice
$\begin{pmatrix} 1 & 2\\ 2 & 1\end{pmatrix}$. Un produit scalaire
est exactement une forme bilinéaire \omterm{def:b2:quadratic:adjoint}{symétrique} dont la \omterm{def:b2:quadratic:def}{forme
quadratique} est définie positive~; ce chapitre étudie le cas
général, de signe indéfini.
\end{example}

\begin{example}[La congruence en action]\label{ex:b2:quadratic:congruence}
Prenons $q(x, y) = x^2 + 4xy + y^2$ (de matrice $B = \begin{pmatrix}
1 & 2\\ 2 & 1\end{pmatrix}$) et la nouvelle base $e_1' = (1, 1)$,
$e_2' = (1, -1)$, c'est-à-dire $P = \begin{pmatrix} 1 & 1\\ 1 &
-1\end{pmatrix}$. Alors
\[
P^{\mathsf T}BP
= \begin{pmatrix} 1 & 1\\ 1 & -1\end{pmatrix}
\begin{pmatrix} 1 & 2\\ 2 & 1\end{pmatrix}
\begin{pmatrix} 1 & 1\\ 1 & -1\end{pmatrix}
= \begin{pmatrix} 6 & 0\\ 0 & -2\end{pmatrix} :
\]
dans les coordonnées $(u, v)$ selon la nouvelle base, $q = 6u^2 -
2v^2$ --- vérification~: $x = u + v$, $y = u - v$ donne directement
$x^2 + 4xy + y^2 = 6u^2 - 2v^2$. Remarquons que les nouveaux
coefficients diagonaux $6, -2$ ne sont \emph{pas} les \omterm{def:b2:reduction:eigen}{valeurs
propres} $3, -1$ de $B$~: la \omterm{def:b2:quadratic:def}{congruence} redimensionne, seule la
similitude préserve les \omterm{def:b2:reduction:eigen}{spectres} --- mais les signes coïncident,
comme l'exige le théorème de Sylvester. (La base est ici orthogonale
mais non orthonormée~; la normaliser par $\frac{1}{\sqrt2}$
diviserait la diagonale par $2$ et redonnerait les \omterm{def:b2:reduction:eigen}{valeurs propres}.)
\end{example}

\begin{example}[Les déterminants de Gram mesurent une aire]\label{ex:b2:quadratic:gramarea}
Pour $v_1, v_2$ dans un espace euclidien, la matrice de Gram $G =
\bigl(\langle v_i, v_j\rangle\bigr)$ encode les longueurs et
l'angle~; son \omterm{def:b2:linalg:det}{déterminant} encode l'\emph{aire}~:
\[
\det G = \norm{v_1}^2\norm{v_2}^2 - \langle v_1, v_2\rangle^2
= \norm{v_1}^2\norm{v_2}^2\bigl(1 - \cos^2\theta\bigr)
= \bigl(\norm{v_1}\,\norm{v_2}\sin\theta\bigr)^2 ,
\]
l'aire au carré du parallélogramme construit sur $v_1, v_2$ ---
et Cauchy--Schwarz n'est autre que l'énoncé $\det G \geq 0$.
Cas concret~: $v_1 = (1, 2, 2)$, $v_2 = (2, 1, -2)$ dans
$\R^3$~:
\[
G = \begin{pmatrix} 9 & 0\\ 0 & 9 \end{pmatrix},
\qquad
\det G = 81 :
\]
les vecteurs sont orthogonaux, de longueur $3$, et engendrent un
parallélogramme (ici un carré) d'aire $\sqrt{81} = 9$.
Éclairage final~: aucun produit vectoriel ni magie propre à la
dimension~$3$ n'a été utilisé --- $\sqrt{\det G}$ mesure le volume
$k$-dimensionnel en \emph{toute} dimension, point de départ de la
partie~I du problème du week-end et des intégrales d'aire de
surface plus loin dans ce volume.
\end{example}

\section{Réduction de Gauss et inertie de Sylvester}

\begin{theorem}[Réduction de Gauss]\label{thm:b2:quadratic:gauss}
Toute \omterm{def:b2:quadratic:def}{forme quadratique} $q$ sur un espace réel de dimension finie
peut s'écrire\index{réduction de Gauss}
\[
q = \sum_{i=1}^{s} \ell_i^2 - \sum_{j=1}^{t} m_j^2 ,
\]
où $\ell_1, \dots, \ell_s, m_1, \dots, m_t$ sont des formes
linéaires linéairement indépendantes~; de manière équivalente, une
certaine base rend la matrice de $q$ diagonale, à coefficients $+1$
($s$ fois), $-1$ ($t$ fois), $0$.
\end{theorem}

\begin{proof}
Récurrence sur le nombre de variables, en coordonnées~: $q(x_1,
\dots, x_n)$.

\emph{Cas 1~: un carré apparaît}, disons que le coefficient $a$ de
$x_1^2$ est non nul. On regroupe tous les termes en $x_1$ et on
complète le carré~:
\[
q = a\Bigl(x_1 + \frac{1}{a}\,\lambda(x_2, \dots,
x_n)\Bigr)^{\!2} + q_1(x_2, \dots, x_n),
\]
où $\lambda$ est linéaire et $q_1$ quadratique en les variables
restantes~: une forme indépendante s'est détachée (elle fait
intervenir $x_1$, les autres non), la récurrence s'applique à
$q_1$, et les signes $\pm$ proviennent du signe de $a$ après
redimensionnement par $\sqrt{\abs a}$.

\emph{Cas 2~: aucun carré, mais un terme croisé}, disons
$b\,x_1x_2$ avec $b \neq 0$. On utilise l'identité
\[
x_1x_2 = \tfrac14\bigl((x_1 + x_2)^2 - (x_1 - x_2)^2\bigr)
\]
après regroupement~: en écrivant $q = b\,x_1x_2 + x_1\alpha +
x_2\beta + q_2$ (avec $\alpha, \beta, q_2$ en les autres variables),
on vérifie
\[
q = \frac{b}{4}\Bigl[\Bigl(x_1 + x_2 + \frac{\alpha +
\beta}{b}\Bigr)^{2} - \Bigl(x_1 - x_2 + \frac{\beta -
\alpha}{b}\Bigr)^{2}\Bigr] + \widetilde q ,
\]
avec $\widetilde q$ ne dépendant pas de $x_1, x_2$~: deux formes
indépendantes se détachent, et la récurrence conclut.

Indépendance des formes collectées~: ordonnons les paquets dans
l'\omterm{def:b2:structures:generated}{ordre} de production. Les formes du premier paquet contiennent
$x_1$ (cas~1) ou $x_1, x_2$ (cas~2)~; toutes les formes suivantes ne
dépendent pas de ces variables. Supposons qu'une combinaison
linéaire de toutes les formes collectées s'annule. En lisant le
coefficient de $x_1$ (et de $x_2$)~: seul le premier paquet
contribue, et à l'intérieur de ce paquet la ou les deux formes sont
manifestement indépendantes ($\ell$ seule~; ou $\ell \pm m$ avec
$\ell, m$ indépendantes)~: les coefficients du premier paquet
s'annulent. On retire le paquet et on recommence~: par récurrence le
long des paquets, tous les coefficients s'annulent --- la famille
tout entière est libre, la triangularité rendue explicite.
\end{proof}

\begin{theorem}[Loi d'inertie de Sylvester]\label{thm:b2:quadratic:sylvester}
Le couple $(s, t)$ du~\cref{thm:b2:quadratic:gauss} ne dépend que de
$q$, pas de la réduction~: c'est la
\emph{signature}\index{signature d'une forme quadratique} de $q$.
De plus
\[
s = \max\{\dim F : q|_F \text{ définie positive}\},
\]
et symétriquement pour $t$.
\end{theorem}

\begin{proof}
Soit $q = \sum_{i \leq s}\ell_i^2 - \sum_{j\leq t} m_j^2$ et soit
$F_+$ le sous-espace engendré par les vecteurs (pré-)duaux sur
lesquels $(\ell_i)$ se restreignent en coordonnées ---
concrètement~: on complète la famille indépendante $(\ell_1, \dots,
\ell_s, m_1, \dots, m_t)$ en une base du \omterm{def:b2:linalg:dual}{dual} $E^*$, et on note
$(u_1, \dots, u_n)$ la base de $E$ dont ce sont les formes
coordonnées (la base pré-duale~: $\ell_i(u_k) = \delta_{ik}$ pour $k
\leq s$, et les formes suivantes s'annulent sur les vecteurs
précédents). Posons $F_+ = \operatorname{Vect}(u_1, \dots, u_s)$~:
pour $x = \sum_{i\leq s}x_iu_i \in F_+$,
\[
\ell_i(x) = x_i,
\qquad m_j(x) = 0,
\qquad\text{donc}\qquad
q(x) = \sum_{i\leq s}x_i^2 > 0 \quad (x \neq 0) :
\]
$q|_{F_+}$ est définie positive et le maximum de la formule est
$\geq s$. Réciproquement, soit $F$ un sous-espace quelconque tel que
$q|_F$ soit définie positive, et $G = \{x : \ell_1(x) = \dots =
\ell_s(x) = 0\}$, de codimension $\leq s$~; sur $G$, $q(x) = -\sum
m_j^2 \leq 0$. Alors $F \cap G = \{0\}$ (un vecteur non nul y
vérifierait $q > 0$ et $q \leq 0$), donc $\dim F \leq \dim E - \dim
G \leq s$. Ainsi le maximum vaut $s$ pour \emph{toute} réduction~:
$s$ est intrinsèque, et $t = \operatorname{rk} q - s$ de même.
\end{proof}

\begin{example}\label{ex:b2:quadratic:gaussexample}
$q(x, y, z) = xy + yz + zx$ (aucun carré). Avec $x_1 = x$, $x_2 =
y$~: $q = xy + z(x + y)$, et l'identité des deux carrés donne
\[
q = \tfrac14(x + y + 2z)^2 - \tfrac14(x - y)^2 - z^2 ,
\]
(développer pour vérifier). Trois formes indépendantes~: signature
$(1, 2)$, rang $3$. Une direction positive, deux négatives~: la
géométrie en «~cône de lumière~» de cette forme.
\end{example}

\begin{example}[Une forme dégénérée, réduite complètement]\label{ex:b2:quadratic:degenerate}
$q(x, y, z) = xy + yz$ sur $\R^3$~: aucun carré, donc cas~2 avec le
regroupement $q = y(x + z)$. L'identité des deux carrés sur le
produit des formes indépendantes $y$ et $x + z$~:
\[
q = \frac14\bigl(y + x + z\bigr)^2 -
\frac14\bigl(y - x - z\bigr)^2 .
\]
Les deux formes linéaires $y + x + z$ et $y - x - z$ sont
indépendantes (leur différence est $2(x+z)$, leur somme $2y$), donc
Sylvester se lit directement~: signature $(1, 1)$, rang $2$ ---
\emph{dégénérée}. Le noyau de la forme polaire s'obtient en
résolvant $\varphi(v, \cdot) = 0$~: avec la matrice
$\frac12\begin{pmatrix} 0&1&0\\ 1&0&1\\ 0&1&0\end{pmatrix}$,
le noyau est $\{y = 0,\ x + z = 0\} = \R\,(1, 0, -1)$, la
direction selon laquelle $q$ ne voit rien. Éclairage final~: le
défaut de rang se manifeste dans Gauss comme un «~manque de
variables~» --- la réduction n'a produit que deux carrés sur trois
dimensions, et la dimension manquante est exactement le noyau.
\end{example}

\begin{example}[Une forme, deux chemins vers la signature]\label{ex:b2:quadratic:tworoads}
$q(x, y, z) = 2x^2 + 2y^2 + 2z^2 + 2xy + 2yz$, de matrice
$\begin{pmatrix} 2 & 1 & 0\\ 1 & 2 & 1\\ 0 & 1 & 2
\end{pmatrix}$. \emph{Chemin~1, Gauss~:} on complète les carrés dans
l'\omterm{def:b2:structures:generated}{ordre},
\[
q = 2\Bigl(x + \frac y2\Bigr)^{\!2} + \frac32 y^2 + 2yz + 2z^2
= 2\Bigl(x + \frac y2\Bigr)^{\!2}
+ \frac32\Bigl(y + \frac{2z}{3}\Bigr)^{\!2} + \frac43 z^2 :
\]
trois carrés positifs sur des formes indépendantes, signature $(3,
0)$~: définie positive. \emph{Chemin~2, \omterm{def:b2:reduction:eigen}{valeurs propres}~:} la
matrice est la tridiagonale $2I + N$ où $N$ est la matrice des
voisins~; ses \omterm{def:b2:reduction:eigen}{valeurs propres} sont $2 + \sqrt2$, $2$, $2 - \sqrt2$
(vérifier les vecteurs propres $(1, \pm\sqrt2, 1)$ et $(1, 0, -1)$),
toutes positives~: même verdict, d'après le~
\cref{cor:b2:quadratic:principalaxes}. Gauss est plus rapide~; les
\omterm{def:b2:reduction:eigen}{valeurs propres} en disent plus (elles donnent les axes principaux et
les valeurs extrêmes de $q$ sur la sphère). Éclairage final~: les
pivots positifs $2, \frac32, \frac43$ de Gauss sont exactement les
rapports $\frac{\Delta_k}{\Delta_{k-1}}$ des mineurs principaux
dominants ($\Delta_1 = 2$, $\Delta_2 = 3$, $\Delta_3 = 4$) --- le
problème du week-end le prouve en toute généralité.
\end{example}

\begin{method}[Calculer une signature~: trois voies]\label{met:b2:quadratic:signature}
\begin{enumerate}
  \item \emph{Gauss} (marche toujours, le plus rapide à la main)~:
        compléter les carrés dans l'\omterm{def:b2:structures:generated}{ordre}, cas~2 lorsqu'aucun carré
        n'est disponible~; compter les signes. Vérifier que les
        formes linéaires collectées sont indépendantes --- moins de
        formes que de variables signale un noyau
        (\cref{ex:b2:quadratic:degenerate}).
  \item \emph{Mineurs dominants} (pour les tests de définitude)~:
        tous les $\Delta_k > 0$ ssi définie positive
        (\cref{exo:b2:quadratic:8})~; les pivots
        $\Delta_k/\Delta_{k-1}$ donnent même les coefficients de
        Gauss (problème du week-end). Échoue silencieusement si un
        $\Delta_k = 0$~: revenir à la voie~1.
  \item \emph{\omterm{def:b2:reduction:eigen}{Valeurs propres}} (le plus informatif, le plus
        coûteux)~: signes du \omterm{def:b2:reduction:eigen}{spectre}
        (\cref{cor:b2:quadratic:principalaxes})~; fournit aussi les
        axes principaux et les valeurs extrêmes de $q$ sur la sphère
        unité. À privilégier lorsque la structure propre est de
        toute façon nécessaire.
\end{enumerate}
\end{method}

\section{Le théorème spectral}

Soit désormais $E$ euclidien (produit scalaire
$\langle\cdot,\cdot\rangle$, volume de l'année~1).

\begin{definition}[Adjoint~; endomorphismes symétriques]\label{def:b2:quadratic:adjoint}
Pour $u \in \mathcal{L}(E)$, l'\emph{adjoint}\index{adjoint}
$u^*$ est l'unique endomorphisme vérifiant
\[
\langle u(x), y\rangle = \langle x, u^*(y)\rangle
\qquad (x, y \in E);
\]
dans une base orthonormée, $\operatorname{Mat}(u^*) =
\operatorname{Mat}(u)^{\mathsf T}$. On dit que $u$ est
\emph{symétrique}\index{endomorphisme symétrique} (auto-adjoint)
lorsque $u^* = u$ --- de manière équivalente, sa matrice dans une
base orthonormée est symétrique.
\end{definition}

\begin{proof}[Existence et unicité de l'adjoint]
Pour $y$ fixé, la forme $x \mapsto \langle u(x), y\rangle$ est
linéaire, donc (dimension finie) de la forme $\langle x,
z_y\rangle$ pour un unique $z_y$ --- l'application $y \mapsto z_y =:
u^*(y)$ étant linéaire par unicité. L'identification matricielle~:
$\langle u(e_i), e_j\rangle$ lu dans les deux sens.
\end{proof}

\begin{example}[L'adjoint dépend du produit scalaire]\label{ex:b2:quadratic:weightedadjoint}
Sur $\R^2$, prenons le produit scalaire \emph{pondéré} $\langle x,
y\rangle_D = x_1y_1 + 2x_2y_2$ (de matrice $D =
\operatorname{diag}(1,2)$) et $u$ de matrice $A =
\begin{pmatrix} 0 & 1\\ 0 & 0\end{pmatrix}$ dans la base
canonique. À partir de $\langle u(x), y\rangle_D = (Ax)^{\mathsf
T}Dy = x^{\mathsf T}(A^{\mathsf T}D)y$ et de $\langle x,
u^*(y)\rangle_D = x^{\mathsf T}(DA^*)y$, la matrice de l'\omterm{def:b2:quadratic:adjoint}{adjoint} est
\[
A^* = D^{-1}A^{\mathsf T}D
= \begin{pmatrix} 1 & 0\\ 0 & \tfrac12\end{pmatrix}
\begin{pmatrix} 0 & 0\\ 1 & 0\end{pmatrix}
\begin{pmatrix} 1 & 0\\ 0 & 2\end{pmatrix}
= \begin{pmatrix} 0 & 0\\ \tfrac12 & 0\end{pmatrix}
\neq A^{\mathsf T} .
\]
Vérification sur $x = (1,0)$, $y = (0,1)$~:
\[
\langle u(x), y\rangle_D = \langle (0,0), y\rangle_D = 0,
\quad
\langle x, u^*(y)\rangle_D
= \langle(1,0), (0,\tfrac12)\rangle_D = 0 ;
\]
sur $x = (0,1)$, $y = (1,0)$~:
\[
\langle u(x), y\rangle_D = \langle(1,0),(1,0)\rangle_D = 1,
\quad
\langle x, u^*(y)\rangle_D =
\langle(0,1),(0,\tfrac12)\rangle_D = 1 .
\] Éclairage final~: «~$\operatorname{Mat}(u^*)
= \operatorname{Mat}(u)^{\mathsf T}$~» est un énoncé sur les bases
\emph{orthonormées} uniquement~; en général la métrique $D$
intervient, exactement comme dans la réduction simultanée du
problème du week-end.
\end{example}

\begin{theorem}[Théorème spectral]\label{thm:b2:quadratic:spectral}
Soit $u$ un \omterm{def:b2:quadratic:adjoint}{endomorphisme symétrique} d'un espace euclidien $E$.
Alors $E$ possède une \emph{base orthonormée de vecteurs propres} de
$u$~; toutes les \omterm{def:b2:reduction:eigen}{valeurs propres} sont réelles, et les sous-espaces
propres associés à des \omterm{def:b2:reduction:eigen}{valeurs propres} distinctes sont orthogonaux.
Forme matricielle~: toute matrice \omterm{def:b2:quadratic:adjoint}{symétrique} réelle $A$ s'écrit
\[
A = P\,D\,P^{\mathsf T},
\qquad P \text{ orthogonale } (P^{\mathsf T} P = I),\ D
\text{ diagonale}.\index{théorème spectral}
\]
\end{theorem}

\begin{proof}
\emph{Existence d'un vecteur propre.} La fonction $x \mapsto \langle
u(x), x\rangle$ est \omterm{def:b2:metric:continuity}{continue} sur la sphère unité $S$ de $E$, qui est
compacte (dimension finie, \cref{thm:b2:nvs:finitedim})~: elle
atteint son maximum $\lambda$ en un point $a \in S$. Montrons que
$u(a) = \lambda a$. Pour tout $y \perp a$ avec $\norm y = 1$ et $t
\in \R$, le vecteur $x_t = \frac{a + ty}{\sqrt{1 + t^2}}$ est sur
$S$ ($\norm{a + ty}^2 = 1 + t^2$ par Pythagore)~; en développant la
fonction maximisée,
\[
g(t) = \langle u(x_t), x_t\rangle
= \frac{\langle u(a), a\rangle + 2t\langle u(a), y\rangle +
t^2\langle u(y), y\rangle}{1 + t^2}
\]
(la symétrie de $u$ a fusionné les deux termes croisés~: $\langle
u(a), y\rangle = \langle a, u(y)\rangle = \langle u(y), a\rangle$).
$g$ est une fonction dérivable de $t$ ayant un maximum en $t = 0$~;
la règle de dérivation d'un quotient en $0$ donne
\[
g'(0) = \frac{2\langle u(a), y\rangle\cdot 1 - \langle u(a),
a\rangle\cdot 0}{1} = 2\langle u(a), y\rangle = 0 .
\]
Ainsi $u(a)$ est orthogonal à l'hyperplan $a^\perp$ tout entier~:
$u(a) \in (a^{\perp})^{\perp} = \R a$, c'est-à-dire $u(a) = \mu a$~;
et $\mu = \langle u(a), a\rangle = \lambda$.

\emph{Récurrence.} Le supplémentaire orthogonal $F = a^\perp$ est
$u$-stable~: pour $x \perp a$, $\langle u(x), a\rangle = \langle x,
u(a)\rangle = \lambda\langle x, a\rangle = 0$. La restriction
$u|_F$ est \omterm{def:b2:quadratic:adjoint}{symétrique} pour le produit scalaire induit~; par
récurrence sur la dimension, $F$ possède une base orthonormée de
vecteurs propres~; on préfixe $a$.

\emph{Compléments.} Les \omterm{def:b2:reduction:eigen}{valeurs propres} sont les nombres réels
$\langle u(e), e\rangle$ sur la base propre. Orthogonalité des
sous-espaces propres~: $u(x) = \lambda x$, $u(y) = \mu y$ donnent
$\lambda\langle x, y\rangle = \langle u(x), y\rangle = \langle x,
u(y)\rangle = \mu\langle x, y\rangle$, donc $\langle x, y\rangle =
0$ lorsque $\lambda \neq \mu$. Forme matricielle~: les colonnes de
$P$ = la base orthonormée de vecteurs propres.
\end{proof}

\begin{example}[Un calcul spectral complet]\label{ex:b2:quadratic:fullrun}
Diagonaliser orthogonalement $A = \begin{pmatrix} 1 & 2\\ 2 &
1\end{pmatrix}$. \omterm{def:b2:reduction:charpoly}{Polynôme caractéristique} $(1 - \lambda)^2 -
4$~: \omterm{def:b2:reduction:eigen}{valeurs propres} $3$ et $-1$. Vecteurs propres~: $(A - 3I)v = 0$
donne $v_1 = \frac{1}{\sqrt2}(1,1)$~; $(A + I)v = 0$ donne $v_2
= \frac{1}{\sqrt2}(1,-1)$ --- orthogonaux, comme le~
\cref{thm:b2:quadratic:spectral} le garantit sans calcul. Avec $P =
(v_1\ v_2)$ (une rotation d'angle $\frac\pi4$)~:
\[
P^{\mathsf T}AP = \begin{pmatrix} 3 & 0\\ 0 & -1
\end{pmatrix},
\qquad
x^2 + 4xy + y^2 = 3u^2 - v^2
\quad\text{dans le repère tourné} .
\]
Ainsi la forme de l'\cref{exo:b2:quadratic:1} est une forme de type
hyperbolique~: signature $(1,1)$, cohérente avec sa \omterm{thm:b2:quadratic:gauss}{réduction de
Gauss} $(x + 2y)^2 - 3y^2$ --- des carrés différents, la même
signature, comme l'exige Sylvester. Éclairage final~: Gauss a donné
la réponse plus vite, mais la voie spectrale signale aussi que sur
le cercle unité $q$ parcourt exactement $\intcc{-1}{3}$, atteint le
long de $v_2$ et $v_1$~: le travail supplémentaire achète de la
géométrie.
\end{example}

\begin{corollary}[Axes principaux~; tests de positivité]\label{cor:b2:quadratic:principalaxes}
\begin{enumerate}
  \item Toute \omterm{def:b2:quadratic:def}{forme quadratique} $q$ sur un espace euclidien se
        diagonalise dans une base \emph{orthonormée}~: $q(x) =
        \sum_i \lambda_i x_i^2$ où les $\lambda_i$ sont les \omterm{def:b2:reduction:eigen}{valeurs
        propres} de la matrice \omterm{def:b2:quadratic:adjoint}{symétrique} de $q$~; la signature
        compte les \omterm{def:b2:reduction:eigen}{valeurs propres} positives et négatives.
  \item Une matrice \omterm{def:b2:quadratic:adjoint}{symétrique} est semi-définie positive (resp.\
        définie positive) ssi toutes ses \omterm{def:b2:reduction:eigen}{valeurs propres} sont $\geq
        0$ (resp.\ $> 0$)~; et alors les valeurs extrêmes du
        quotient de Rayleigh sont
        \[
        \min_{\norm x = 1} \langle Ax, x\rangle = \lambda_{\min},
        \qquad
        \max_{\norm x = 1} \langle Ax, x\rangle = \lambda_{\max} .
        \]
\end{enumerate}
\end{corollary}

\begin{proof}
(1) Écrivons $q(x) = \langle A x, x\rangle$ avec $A$ \omterm{def:b2:quadratic:adjoint}{symétrique} (la
matrice de $q$ dans une base orthonormée)~; diagonalisons $A$ par le
théorème spectral~: pour $x = \sum x_ie_i$ dans la base propre
orthonormée,
\[
q(x) = \Bigl\langle \sum_i \lambda_ix_ie_i,\ \sum_j
x_je_j\Bigr\rangle = \sum_i \lambda_i x_i^2
\]
(l'orthonormalité tue les termes croisés). Les signes $\pm$ des
$\lambda_i$ comptent la signature d'après Sylvester~: redimensionner
chaque coordonnée par $\sqrt{\abs{\lambda_i}}$ exhibe une \omterm{thm:b2:quadratic:gauss}{réduction
de Gauss} à formes indépendantes.

(2) Dans la base propre, $\langle Ax, x\rangle = \sum \lambda_i
x_i^2$, encadré entre $\lambda_{\min}\norm x^2$ et
$\lambda_{\max}\norm x^2$, avec égalité aux vecteurs propres
correspondants~; la positivité de toutes les \omterm{def:b2:reduction:eigen}{valeurs propres}
équivaut donc à la positivité de la forme.
\end{proof}

\begin{example}\label{ex:b2:quadratic:spectralexample}
$A = \begin{pmatrix} 2 & 1\\ 1 & 2 \end{pmatrix}$~: \omterm{def:b2:reduction:eigen}{valeurs propres}
$3$ (vecteur propre $\frac{1}{\sqrt2}(1,1)$) et $1$
($\frac{1}{\sqrt2}(1,-1)$). La \omterm{def:b2:quadratic:def}{forme quadratique} $2x^2 + 2xy + 2y^2$
devient $3X^2 + Y^2$ dans le repère orthonormé tourné~: les axes
principaux d'une ellipse, calculés. La \omterm{thm:b2:quadratic:gauss}{réduction de Gauss} atteint
elle aussi une forme diagonale, mais seul le théorème spectral
l'atteint \emph{sans distordre les longueurs}.
\end{example}

\begin{example}[Une ellipse entièrement identifiée]\label{ex:b2:quadratic:ellipserun}
Quelle courbe est $5x^2 + 4xy + 2y^2 = 6$~? La matrice
$\begin{pmatrix} 5 & 2\\ 2 & 2\end{pmatrix}$ a pour \omterm{def:b2:reduction:charpoly}{polynôme
caractéristique} $\lambda^2 - 7\lambda + 6 = (\lambda - 1)(\lambda -
6)$~: \omterm{def:b2:reduction:eigen}{valeurs propres} $1$ et $6$, toutes deux positives --- une
ellipse. Vecteurs propres orthonormés~: pour $\lambda = 1$, on
résout $\begin{pmatrix} 4 & 2\\ 2 & 1\end{pmatrix}v = 0$~: $v_1 =
\frac{1}{\sqrt5}(1, -2)$~; pour $\lambda = 6$~: $v_2 =
\frac{1}{\sqrt5}(2, 1)$. Dans les coordonnées tournées $(X, Y)$
selon $(v_2, v_1)$ l'équation devient
\[
6X^2 + Y^2 = 6,
\qquad\text{c.-à-d.}\qquad
X^2 + \frac{Y^2}{6} = 1 :
\]
demi-axes $1$ (selon $v_2$) et $\sqrt6$ (selon $v_1$).
Éclairage final~: l'allure grossière était gratuite --- $\det
= 6 > 0$ et une trace positive annoncent une ellipse avant tout
calcul de vecteur propre --- mais seul le théorème spectral livre
les directions et longueurs des axes, c'est-à-dire la géométrie
réelle.
\end{example}

\begin{example}[Extrema sur la sphère, lus sur le spectre]\label{ex:b2:quadratic:sphereextremes}
Quelles sont les valeurs extrêmes de $q(x,y,z) = 2xy + 2yz + 2zx$
sur la sphère unité~? Sa matrice (celle dont la partie
hors-diagonale ne contient que des uns, de l'\cref{exo:b2:quadratic:2})
a pour \omterm{def:b2:reduction:eigen}{valeurs propres} $2$ et $-1$ (double), donc d'après le~
\cref{cor:b2:quadratic:principalaxes}~(2)~:
\[
\max_{\norm v = 1} q(v) = 2
\ \text{ en } v = \tfrac{1}{\sqrt3}(1,1,1),
\qquad
\min_{\norm v = 1} q(v) = -1
\ \text{ sur le cercle } x + y + z = 0 .
\]
Aucun calcul différentiel, aucun multiplicateur de Lagrange~: le
théorème spectral résout d'emblée cette optimisation sous contrainte
--- et exhibe le maximiseur. Éclairage final~: comparer avec la
méthode des multiplicateurs du chapitre de calcul différentiel, qui
trouve les mêmes points critiques au prix de plus de travail~; pour
des objectifs \emph{quadratiques} sur des sphères, les \omterm{def:b2:reduction:eigen}{spectres} sont
la voie royale (le problème du week-end du chapitre hermitien
construit toute la théorie de Courant--Fischer sur cette
observation).
\end{example}

\begin{remark}[Pièges classiques]
\emph{(i) La \omterm{def:b2:quadratic:def}{congruence} n'est pas la similitude~:} un changement de
base pour une forme agit par $P^{\mathsf T}BP$, non par $P^{-1}BP$~;
les \omterm{def:b2:reduction:eigen}{valeurs propres} ne sont \emph{pas} des invariants d'une \omterm{def:b2:quadratic:def}{forme
quadratique} ($I$ et $4I$ sont congruentes via $P = 2I$) --- seuls
leurs signes le sont (Sylvester). Ne parler des \omterm{def:b2:reduction:eigen}{valeurs propres}
d'une forme qu'une fois un produit scalaire fixé. \emph{(ii) Des
coefficients positifs ne prouvent rien~:} $\begin{pmatrix} 1 & 2\\ 2
& 1\end{pmatrix}$ a tous ses coefficients positifs et pourtant une
signature $(1,1)$ ($\det = -3$)~; réciproquement une matrice définie
positive peut avoir des coefficients hors-diagonaux négatifs
(l'\cref{ex:b2:quadratic:tworoads} décalé~: $2I - N$ convient tout
autant). Utiliser la~\cref{met:b2:quadratic:signature}. \emph{(iii)
Carrés dépendants~:} écrire $q = \ell_1^2 - \ell_2^2$ ne dit rien si
$\ell_1, \ell_2$ sont proportionnelles --- $x^2 + 2xy + y^2 =
(x+y)^2$ est de rang $1$, non $2$~; toujours vérifier
l'indépendance avant de lire la signature. \emph{(iv) Extrema sur la
sphère sans compacité~:} les bornes de Rayleigh du~
\cref{cor:b2:quadratic:principalaxes} sont atteintes parce que la
sphère est compacte~; sur la boule ouverte ou l'espace entier, une
forme indéfinie n'a ni maximum ni minimum.
\end{remark}

\begin{remark}[Où cela sert]
Le théorème spectral est le résultat le plus exporté de ce livre~:
la statistique s'en sert pour diagonaliser les matrices de
covariance (analyse en composantes principales), la mécanique en
extrait les modes propres d'oscillation (la réduction simultanée du
problème du week-end), l'analyse numérique y bâtit les
décompositions de Cholesky et en valeurs singulières (le même
problème), et le chapitre suivant le transporte aux espaces
hermitiens complexes. Le volume de l'année~3 en prouve l'avatar en
dimension infinie pour les opérateurs \omterm{def:b2:quadratic:adjoint}{auto-adjoints} \omterm{def:b2:metric:compact}{compacts}, où
l'argument de compacité de la preuve en dimension finie devient
toute l'histoire.
\end{remark}

\begin{remark}[Perspectives au sein de ce volume]
Les \omterm{def:b2:quadratic:def}{formes quadratiques} traversent le reste du livre~4 sous trois
déguisements. En tant que \emph{hessiennes}~: le chapitre de calcul
différentiel classe les points critiques par la signature de la
forme du second \omterm{def:b2:structures:generated}{ordre}, si bien que l'invariance de Sylvester est ce
qui donne un sens précis au mot «~selle~». En tant
qu'\emph{énergies}~: les oscillateurs du chapitre d'équations
différentielles portent l'énergie quadratique $\frac12x'^2 +
\frac12\omega^2x^2$, et la réduction simultanée du problème du
week-end de ce chapitre est exactement l'extraction des modes
propres. En tant que \emph{géométrie}~: les coniques de ce chapitre
grandissent en les surfaces quadriques des chapitres de géométrie,
où la seconde forme fondamentale d'une surface --- une \omterm{def:b2:quadratic:def}{forme
quadratique} sur chaque plan tangent --- a une signature qui décide
si la surface courbe comme un bol ou comme une selle. Le chapitre
hermitien, ensuite, rejoue toute la partition sur $\C$.
\end{remark}

\section{Exercices}

\begin{exercise}[$\star$]\label{exo:b2:quadratic:1}
Réduire par Gauss et donner le rang et la signature~:
\[
q_1(x,y) = x^2 + 4xy + y^2,
\qquad
q_2(x,y,z) = x^2 + 2y^2 + 3z^2 + 2xy + 2yz .
\]
\end{exercise}

\begin{exercise}[$\star$]\label{exo:b2:quadratic:2}
Diagonaliser orthogonalement $A = \begin{pmatrix} 0 & 1 & 1\\ 1 & 0 &
1\\ 1 & 1 & 0\end{pmatrix}$ (\omterm{def:b2:reduction:eigen}{valeurs propres} issues du calcul du~
\cref{ch:b2:reduction}~; rendre maintenant la base orthonormée) et
réduire la forme $q(x,y,z) = 2xy + 2yz + 2zx$ à ses axes principaux.
\end{exercise}

\begin{exercise}[$\star$]\label{exo:b2:quadratic:3}
Prouver que $u^{**} = u$, $(u \circ v)^* = v^* \circ u^*$, et que
$\ker u^* = (\operatorname{im} u)^{\perp}$. En déduire
$\operatorname{rk} u^* = \operatorname{rk} u$.
\end{exercise}

\begin{exercise}[$\star\star$]\label{exo:b2:quadratic:4}
Soit $A$ \omterm{def:b2:quadratic:adjoint}{symétrique} réelle telle que $A^3 = A$. Prouver que $A^2$
est la matrice d'une projection orthogonale. Plus généralement,
relier les décompositions spectrales de $A$ et de $P(A)$ pour un
polynôme $P$.
\end{exercise}

\begin{exercise}[$\star\star$]\label{exo:b2:quadratic:5}
Prouver que $O(n) = \{P : P^{\mathsf T}P = I\}$ est une partie
compacte de $\mathcal{M}_n(\R)$ \emph{(fermée~: image réciproque de
$\{I\}$ par une application \omterm{def:b2:metric:continuity}{continue}~; bornée~: les colonnes sont
des vecteurs unitaires)}. Est-elle \omterm{def:b2:metric:connected}{connexe}~?
\end{exercise}

\begin{exercise}[$\star\star$]\label{exo:b2:quadratic:6}
(Racine carrée) Soit $A$ \omterm{def:b2:quadratic:adjoint}{symétrique} semi-définie positive.
Construire une matrice \omterm{def:b2:quadratic:adjoint}{symétrique} semi-définie positive $B$ telle
que $B^2 = A$, et prouver qu'elle est \emph{unique} \emph{(existence~:
prendre les racines carrées des \omterm{def:b2:reduction:eigen}{valeurs propres} dans une base
spectrale~; unicité~: un candidat $B$ commute avec $A = B^2$, donc
préserve ses sous-espaces propres --- se ramener au cas scalaire sur
chacun)}.
\end{exercise}

\begin{exercise}[$\star\star$]\label{exo:b2:quadratic:7}
Pour $A$ \omterm{def:b2:quadratic:adjoint}{symétrique} réelle, prouver que $\vertiii{A}_2 :=
\sup_{\norm x_2 = 1}\norm{Ax}_2 = \max_i \abs{\lambda_i}$
(rayon spectral), et calculer $\vertiii{A}_2$ pour $A =
\begin{pmatrix} 1 & 2\\ 2 & 1\end{pmatrix}$.
\end{exercise}

\begin{exercise}[$\star\star\star$]\label{exo:b2:quadratic:8}
(Critère de Sylvester) Soit $A$ \omterm{def:b2:quadratic:adjoint}{symétrique} réelle de mineurs
principaux dominants $\Delta_1, \dots, \Delta_n$ (\omterm{def:b2:linalg:det}{déterminants} des
blocs en haut à gauche). Prouver que $A$ est définie positive si et
seulement si tous les $\Delta_k > 0$. \emph{(Pour $\Rightarrow$~: les
restrictions d'une forme définie sont définies, et le \omterm{def:b2:linalg:det}{déterminant}
d'une matrice définie positive --- le produit de ses \omterm{def:b2:reduction:eigen}{valeurs propres}
--- est positif. Pour $\Leftarrow$~: récurrence sur $n$~; le bloc
$(n-1)$ en haut à gauche est défini positif, diagonaliser la forme
sur ce sous-espace et compléter le carré en la dernière variable~;
le signe du dernier coefficient diagonal est gouverné par $\det A =
\Delta_n > 0$.)}
\end{exercise}

\begin{exercise}[$\star\star\star$]\label{exo:b2:quadratic:9}
(Courant--Fischer, deuxième \omterm{def:b2:reduction:eigen}{valeur propre}) Soit $u$ \omterm{def:b2:quadratic:adjoint}{symétrique} de
\omterm{def:b2:reduction:eigen}{valeurs propres} $\lambda_1 \geq \lambda_2 \geq \dots \geq
\lambda_n$. Prouver
\[
\lambda_2 = \min_{\substack{H \text{ hyperplan}}}\;
\max_{\substack{x \in H,\ \norm x = 1}} \langle u(x), x\rangle .
\]
\emph{(Pour $\leq$~: tout hyperplan rencontre le $2$-plan engendré
par les deux premiers vecteurs propres. Pour $\geq$~: choisir $H =
(e_1)^{\perp}$.)}
\end{exercise}

\begin{exercise}[$\star\star$]\label{exo:b2:quadratic:10}
Déterminer le rang et la signature de $q(x_1, \dots, x_n) = \sum_{i
< j} x_ix_j$ sur $\R^n$ ($n \geq 2$), de deux façons~: par
l'identité algébrique $2q = \bigl(\sum x_i\bigr)^2 - \sum x_i^2$
jointe à la restriction de $q$ à l'hyperplan $\sum x_i = 0$~; et en
calculant les \omterm{def:b2:reduction:eigen}{valeurs propres} de sa matrice $\frac12(J - I)$, où $J$
est la matrice ne contenant que des uns.
\end{exercise}

\begin{exercise}[$\star\star$]\label{exo:b2:quadratic:11}
Soient $A$, $B$ \omterm{def:b2:quadratic:adjoint}{symétriques} réelles avec $B$ semi-définie positive.
Prouver
\[
\lambda_{\min}(A)\operatorname{tr} B
\;\leq\;
\operatorname{tr}(AB)
\;\leq\;
\lambda_{\max}(A)\operatorname{tr} B .
\]
\emph{(Écrire $B = C^{\mathsf T}C$ et $\operatorname{tr}(AB) =
\sum_i \langle A c_i, c_i\rangle$ sur les colonnes $c_i$ de
$C^{\mathsf T}$.)} En particulier $\operatorname{tr}(AB) \geq 0$
lorsque les deux sont semi-définies positives.
\end{exercise}

\begin{exercise}[$\star\star\star$]\label{exo:b2:quadratic:12}
Sur $E = \mathcal{M}_n(\R)$, on considère $q(M) =
\operatorname{tr}(M^2)$.
\begin{enumerate}
  \item Montrer que $q$ est une \omterm{def:b2:quadratic:def}{forme quadratique} de forme polaire
        $\varphi(M, N) = \operatorname{tr}(MN)$.
  \item Montrer que les matrices \omterm{def:b2:quadratic:adjoint}{symétriques} et antisymétriques
        forment des sous-espaces $\varphi$-orthogonaux sur lesquels
        $q$ est respectivement définie positive et définie négative
        \emph{(calculer $\operatorname{tr}(M^2)$ coefficient par
        coefficient dans chaque cas)}.
  \item Conclure~: $q$ a pour signature $\bigl(\frac{n(n+1)}{2},
        \frac{n(n-1)}{2}\bigr)$ et rang $n^2$.
\end{enumerate}
\end{exercise}

\section{Problème~: Cholesky, Hadamard et la décomposition
polaire}

\begin{problem}\label{pb:b2:quadratic:1}
Le théorème spectral est un microscope~; ce problème l'utilise comme
une usine. À partir des matrices de Gram, on fabrique la
\emph{factorisation de Cholesky}\index{factorisation de Cholesky}
(et on identifie les pivots de Gauss aux rapports de mineurs), puis
on prouve l'\emph{inégalité de Hadamard}\index{inégalité de
Hadamard} sur les \omterm{def:b2:linalg:det}{déterminants}, on construit la \emph{décomposition
polaire}\index{décomposition polaire} $A = QS$ et la
\emph{décomposition en valeurs singulières}, on classe les coniques
planes, et on termine par la réduction simultanée de deux formes ---
le théorème derrière les modes propres d'oscillation. Dans tout ce
qui suit, $E = \R^n$ muni de son produit scalaire canonique.

\medskip
\noindent\textbf{Partie~I --- Matrices de Gram et Cholesky.} Pour
des vecteurs $v_1, \dots, v_n \in E$, leur \emph{matrice de Gram}
est $G = \bigl(\langle v_i, v_j\rangle\bigr)_{i,j}$.
\begin{enumerate}
  \item Montrer que $G$ est \omterm{def:b2:quadratic:adjoint}{symétrique} semi-définie positive, et
        définie positive si et seulement si $(v_1, \dots, v_n)$ est
        linéairement indépendante \emph{(calculer $X^{\mathsf T}GX$)}.
  \item Réciproquement, montrer que toute matrice \omterm{def:b2:quadratic:adjoint}{symétrique}
        semi-définie positive $A$ est une matrice de Gram~: $A =
        C^{\mathsf T}C$ pour un certain $C$ (utiliser la racine
        carrée de l'\cref{exo:b2:quadratic:6}), avec $C$ inversible
        ssi $A$ est définie.
  \item En déduire qu'une matrice semi-définie positive $A$ vérifie
        $\abs{a_{ij}} \leq \sqrt{a_{ii}\,a_{jj}}$ pour tous $i,
        j$ \emph{(restreindre à deux coordonnées)} --- l'inégalité
        de Cauchy--Schwarz, relue matriciellement.
  \item (Cholesky) Soit $A$ définie positive. Prouver qu'il existe
        une \emph{unique} matrice triangulaire supérieure $T$ à
        coefficients diagonaux positifs telle que
        \[
        A = T^{\mathsf T}\,T
        \]
        \emph{(existence~: appliquer Gram--Schmidt à des vecteurs
        réalisant $A$ comme matrice de Gram~; unicité~: si
        $T_1^{\mathsf T}T_1 = T_2^{\mathsf T}T_2$, montrer que
        $T_1T_2^{-1}$ est orthogonale et triangulaire à diagonale
        positive, donc $I$)}.
  \item Montrer que les mineurs principaux dominants vérifient
        $\Delta_k = (t_{11}\cdots t_{kk})^2$, et en déduire que
        les pivots produits par la \omterm{thm:b2:quadratic:gauss}{réduction de Gauss} d'une forme
        définie positive, dans l'\omterm{def:b2:structures:generated}{ordre} naturel des variables, sont
        \[
        d_k = \frac{\Delta_k}{\Delta_{k-1}}
        \qquad (\Delta_0 = 1) :
        \]
        les mineurs du critère de Sylvester
        (\cref{exo:b2:quadratic:8}) et les pivots de Gauss sont
        la même donnée. Vérifier sur
        l'\cref{ex:b2:quadratic:tworoads}.
\end{enumerate}

\medskip
\noindent\textbf{Partie~II --- \omterm{pb:b2:quadratic:1}{Inégalité de Hadamard}.}
\begin{enumerate}[resume]
  \item Soit $A$ définie positive. Prouver
        \[
        \det A \leq a_{11}\,a_{22}\cdots a_{nn}
        \]
        \emph{(normaliser~: $B = DAD$ avec $D =
        \operatorname{diag}(a_{ii}^{-1/2})$ a une diagonale unité~;
        majorer $\det B = \prod \mu_i$ par l'inégalité
        arithmético-géométrique face à $\operatorname{tr} B = n$)}.
  \item Montrer que l'égalité a lieu ssi $A$ est diagonale.
  \item En déduire l'\emph{\omterm{pb:b2:quadratic:1}{inégalité de Hadamard}}~: pour toute
        matrice carrée réelle $M$ de colonnes $c_1, \dots, c_n$,
        \[
        \abs{\det M} \leq \prod_{i=1}^{n}\norm{c_i}_2 ,
        \]
        avec égalité (pour $M$ inversible) ssi les colonnes sont
        deux à deux orthogonales \emph{(appliquer les questions~6--7
        à $M^{\mathsf T}M$)}.
  \item Dividendes géométriques et combinatoires~: interpréter la
        question~8 comme «~le volume d'un parallélépipède est au
        plus le produit de ses longueurs d'arêtes~»~; et montrer
        qu'une matrice dont tous les coefficients sont dans
        $\intcc{-1}{1}$ vérifie $\abs{\det M} \leq n^{n/2}$. (Les
        matrices atteignant cette borne --- les matrices de Hadamard
        --- existent pour $n = 1, 2$ et de nombreux multiples de
        $4$~; savoir si c'est le cas pour \emph{tous} les multiples
        de $4$ est un célèbre problème \omterm{def:b2:metric:topology}{ouvert}.)
\end{enumerate}

\medskip
\noindent\textbf{Partie~III --- \omterm{pb:b2:quadratic:1}{Décomposition polaire} et valeurs
singulières.}
\begin{enumerate}[resume]
  \item Soit $A$ inversible. Montrer que $A^{\mathsf T}A$ est
        définie positive, et que
        \[
        S = \sqrt{A^{\mathsf T}A}
        \quad\text{(la racine carrée de l'\cref{exo:b2:quadratic:6})},
        \qquad Q = AS^{-1}
        \]
        donnent une factorisation $A = QS$ avec $Q$ orthogonale et
        $S$ définie positive.
  \item Prouver que cette factorisation d'un $A$ inversible est
        unique.
  \item Étendre l'existence à un $A$ quelconque~: choisir
        $\varepsilon_k \to 0$ tel que $A + \varepsilon_k I$ soit
        inversible, écrire $A + \varepsilon_kI = Q_kS_k$, et
        utiliser la compacité de $O(n)$
        (\cref{exo:b2:quadratic:5}) pour extraire $Q_k \to Q$~;
        montrer que $S_k = Q_k^{\mathsf T}(A + \varepsilon_kI)$
        \omterm{def:b2:integration:improper}{converge} vers une matrice semi-définie positive $S$ avec $A
        = QS$ et $S = \sqrt{A^{\mathsf T}A}$. Où l'unicité
        échoue-t-elle pour $A$ singulière~?
  \item (Décomposition en valeurs singulières) En déduire que toute
        matrice carrée réelle $A$ s'écrit
        \[
        A = U\,\Sigma\,V^{\mathsf T},
        \qquad U, V \in O(n),\quad
        \Sigma = \operatorname{diag}(\sigma_1, \dots,
        \sigma_n),\ \sigma_i \geq 0 ,
        \]
        où les $\sigma_i$ (les \emph{valeurs singulières}) sont
        les \omterm{def:b2:reduction:eigen}{valeurs propres} de $\sqrt{A^{\mathsf T}A}$.
  \item Trois conséquences~: $\vertiii{A}_2 = \sigma_{\max}$ pour
        \emph{toute} matrice réelle $A$ (généralisant
        l'\cref{exo:b2:quadratic:7})~; $\abs{\det A} =
        \sigma_1\cdots\sigma_n$~; et l'image de la sphère unité par
        un $A$ inversible est un ellipsoïde de demi-axes
        $\sigma_1, \dots, \sigma_n$ le long des colonnes de $U$.
\end{enumerate}

\medskip
\noindent\textbf{Partie~IV --- Coniques, par le théorème spectral.}
Une conique plane est l'ensemble des zéros de $f(x) = q(x) + \langle
b, x\rangle + c$, avec $q \neq 0$ une \omterm{def:b2:quadratic:def}{forme quadratique} de matrice
$A$, $b \in \R^2$, $c \in \R$.
\begin{enumerate}[resume]
  \item Réduire $f$ par une rotation (axes principaux,
        \cref{cor:b2:quadratic:principalaxes}) suivie d'une
        translation, et classer les formes possibles non vides et
        non dégénérées selon la signature de $q$~: ellipse
        ($\det A > 0$), hyperbole ($\det A < 0$), parabole
        ($\det A = 0$, rang $1$, avec le terme linéaire non
        absorbé).
  \item Mener la réduction complète pour
        \[
        x^2 + 4xy + y^2 + 2x - 2y = 4 :
        \]
        coordonnées tournées, équation réduite, nature et
        centre de la conique.
  \item (Coniques à centre) Supposons $\det A \neq 0$. Montrer que
        le \emph{centre} est $x_0 = -\frac12 A^{-1}b$, et que la
        \omterm{def:b2:quadratic:def}{congruence} par $\begin{pmatrix} I & x_0\\ 0 &
        1\end{pmatrix}$ de la matrice $3\times3$ $\widetilde Q
        = \begin{pmatrix} A & b/2 \\ b^{\mathsf T}/2 &
        c\end{pmatrix}$ donne
        \[
        \det\widetilde Q = f(x_0)\,\det A :
        \]
        la conique à centre est dégénérée (un point ou deux droites)
        exactement lorsque $\det\widetilde Q = 0$.
  \item Vérifier la question~17 sur l'exemple de la question~16~:
        calculer $x_0$, $f(x_0)$ et $\det\widetilde Q$, et
        conclure de nouveau que la conique est une hyperbole non
        dégénérée.
  \item (Un faisceau de quadriques) Pour $\lambda \in \R$, classer
        la surface
        \[
        x^2 + y^2 + z^2 + 2\lambda(xy + yz + zx) = 1
        \]
        par les \omterm{def:b2:reduction:eigen}{valeurs propres} de sa matrice \emph{(structure ne
        contenant que des uns~: \omterm{def:b2:reduction:eigen}{valeurs propres} $1 + 2\lambda$ et
        $1 - \lambda$ double)}~: sphère/ellipsoïde, cylindre, paire
        de plans, hyperboloïdes à une et à deux nappes, selon
        $\lambda$.
\end{enumerate}

\medskip
\noindent\textbf{Partie~V --- Deux formes à la fois~: réduction
simultanée.}
\begin{enumerate}[resume]
  \item Soit $q$ définie positive et $q'$ une \omterm{def:b2:quadratic:def}{forme quadratique}
        quelconque sur $E$. Prouver qu'il existe une base de
        $E$ orthonormée pour $q$ et orthogonale pour $q'$~: dans
        celle-ci, $q = \sum x_i^2$ et $q' = \sum \mu_i x_i^2$
        \emph{(utiliser $q$ comme produit scalaire et appliquer le
        théorème spectral à l'endomorphisme représentant $q'$)}.
  \item Forme matricielle~: pour $A$ définie positive et $B$
        \omterm{def:b2:quadratic:adjoint}{symétrique}, il existe un $P$ inversible avec
        $P^{\mathsf T}AP = I$ et $P^{\mathsf T}BP =
        \operatorname{diag}(\mu_1, \dots, \mu_n)$, où les
        $\mu_i$ sont les racines de $\det(B - \mu A) = 0$.
  \item La mener complètement pour
        \[
        A = \begin{pmatrix} 2 & 1\\ 1 & 1 \end{pmatrix},
        \qquad
        B = \begin{pmatrix} 0 & 1\\ 1 & 0 \end{pmatrix} :
        \]
        les \omterm{def:b2:reduction:eigen}{valeurs propres} généralisées $\mu_\pm$, et des vecteurs
        diagonalisant les deux formes.
  \item Montrer que le caractère défini positif ne peut être
        abandonné~: pour
        \[
        A = \begin{pmatrix} 1 & 0\\ 0 & -1 \end{pmatrix},
        \qquad
        B = \begin{pmatrix} 0 & 1\\ 1 & 0 \end{pmatrix},
        \]
        aucune base ne diagonalise les deux formes \emph{(si $P$
        diagonalisait les deux, $\det(B - \mu A)$ serait scindé à
        racines réelles~; le calculer)}.
  \item Montrer que les $\mu_i$ de la question~21 sont les
        \omterm{def:b2:reduction:eigen}{valeurs propres} de $A^{-1}B$, et que $A^{-1}B$, bien que
        non \omterm{def:b2:quadratic:adjoint}{symétrique} en général, est toujours \omterm{def:b2:reduction:diag}{diagonalisable} à
        \omterm{def:b2:reduction:eigen}{valeurs propres} réelles \emph{(conjuguer par $\sqrt A$)}.
  \item Synthèse. En une phrase chacune~: (i) l'unique théorème
        sur lequel chaque partie s'est appuyée~; (ii) quels
        résultats des parties~I--III survivent pour les matrices
        \emph{semi}-définies positives, et lesquels exigent la
        définitude~; (iii) le système physique dont les petites
        oscillations des questions~20--22 diagonalisent (énergie
        cinétique et énergie potentielle comme les deux formes),
        et ce que les $\mu_i$ y signifient.
\end{enumerate}
\end{problem}
